\subsection{Unrelated machines: the same question, a different answer}
\label{sec:unrelated}

Identical machines are a favourable case for exact methods and for saturating the search. To see whether the decomposition is a property of the problem or of the method, the protocol was repeated on
$\mathrm{R}||C_{\max}$, with the same agent, state and decision rules. The IG removes $d$ jobs with one of four strategies and reinserts them greedily, with or without a local search (a descent that moves or
swaps jobs from the critical machine); the ILS applies a best-improvement move, swap or double move from the critical machine. Following the lesson of Section~\ref{sec:f3f4-ablation}, $d/n$ was chosen per size on
separate instances. The synthetic study has $400$ instances (four distributions, five sizes from $50\times10$ to $500\times30$, $T=1{,}000$), with the LP relaxation and a MILP (HiGHS) as references on $160$ of them.

\paragraph{The mechanism matters, with a different cause}
The IG with local search beats the ILS by $4.16\%$ ($384$ of $400$ instances), also at equal CPU time ($2.16\%$). But the \emph{pure} IG does not beat the ILS ($+1.39\%$) and the local search improves it by
$5.40\%$: for unrelated machines the active ingredient is the \emph{combination} of destruction--reconstruction with a descent, whereas for identical machines the cycle alone sufficed.

\paragraph{The policy matters more, but still little, and follows the headroom}
Q-Learning over the four destruction strategies of the IG with local search gains $0.145\%$ over random selection (90\% interval $0.054$--$0.234$), thirteen times the identical-machine value
but thirty times smaller than the mechanism effect; in the ILS the agent is \emph{worse} than random operator selection ($+0.96\%$, $p_{\mathrm{Holm}}<10^{-14}$), also with eight times the budget. The headroom explains the difference:
the omniscient per-instance choice among the strategies gains $0.38\%$ (against $0.04\%$ for identical machines), and the policy gain is large where the headroom is large and null in $U[100,120]$, where the machines
are almost identical. The law was tested on $400$ new instances: the Spearman correlation between the headroom of a cell and the policy gain is $0.59$ over $20$ cells ($p_{\mathrm{Holm}}=0.015$), but the magnitude is unstable
(pooled $0.073\%$, not significant alone; both samples together $0.109\%$, computed after seeing both), and the gain requires the state and temporal-difference learning together (Supplementary Section~\ref{S-sec:m-unrelated}).

\paragraph{Exact references: the metaheuristic does not win here}
The MILP proves optimality in $66.9\%$ of the $160$ instances within $60$ seconds; the IG with local search and Q-Learning ends $2.04\%$ above the best lower bound (MILP: $0.43\%$). At $1000\times50$ ($32$ new
instances) the MILP returns an incumbent in all of them, $2.24\%$ better than the IG at $T=1{,}000$ ($p_{\mathrm{Holm}}=10^{-4}$). On these $32$ instances the policy gain over random selection was $0.52\%$, five times larger than at
$n\leq500$, so a further $96$ instances tested whether it survives, also at equal CPU time (the agent costs $1.39\times$): pooled, it does not ($0.35\%$ and $0.19\%$; $p_{\mathrm{Holm}}=0.19$ and $0.67$; the prediction recorded beforehand was wrong).
The effect is not uniform: in $U[1,100]$ one strategy dominates the mixture by $1.66\%$ and Q-Learning finds it ($+1.73\%$ gain, $17$ wins and no losses), but a strategy validated offline captures the headroom of the fixed strategies at a lower CPU time, and in two
distributions the agent is worse than random selection. The heuristics of this part are therefore not competitive with the exact method at these sizes; their role here is the attribution (Supplementary Section~\ref{S-sec:m-unrelated}).

\subsubsection*{External validation on the standard benchmark}
The synthetic instances were repeated on the standard $\mathrm{R}||C_{\max}$ benchmark of \citet{FanjulPeyroRuiz2010} (the copy distributed with \citealt{AbladEtAl2021}): $1{,}400$ instances (seven types,
$m\in\{10,\ldots,50\}$, $n\in\{100,\ldots,1000\}$) with the best solution known from CPLEX (two hours), Gurobi and an exact method (300 seconds each), proved optimal for $73\%$ of the test instances. Two instances per cell
($280$) were used for tuning, five ($700$) for a pre-specified test (Table~\ref{tab:f31_rcmax_benchmark}) and the rest ($420$) were reserved. The decomposition replicates. The IG with local search beats the ILS at equal CPU
time by $2.14\%$ ($529$ wins against $99$); the policy gain over random selection is $0.055\%$ at equal $T$ and $0.033\%$ at equal CPU (neither confirmed); the validated strategy beats Q-Learning by $0.066\%$
($p_{\mathrm{Holm}}=0.007$); and the headroom law is confirmed on third-party data (Spearman $0.74$ over $28$ cells, $p_{\mathrm{Holm}}=10^{-5}$). The gain is concentrated where one fixed strategy dominates
($0.61\%$ in $U[1,100]$; $+0.20\%$, worse, in machine-correlated instances). The IG with local search ends $2.08\%$ above the best known solution on average, against $0.07\%$ for Gurobi: the exact references remain ahead.

% Auto-generated by analysis/make_tables_cycle5.py (Entrada 129)
\begin{table}[t]
\centering
\caption{External validation on the standard $\mathrm{R}||C_{\max}$ benchmark of \citet{FanjulPeyroRuiz2010} (700 test instances: 7 types $\times$ 5 machine counts $\times$ 4 job counts $\times$ 5 instances; tuning on 280 other instances; $T=1000$). $\bar d$ in \% of the makespan (negative: first better); Wilcoxon, Holm over five. B3 is the Spearman correlation between the headroom of the destruction strategies and the policy gain over the 28 cells (type $\times$ $n$), one-sided. Below the line, the mean GAP (\%) to the best known solution (the minimum of CPLEX in 2 h, Gurobi and the LR method in 300 s; \citealt{AbladEtAl2021}), by instance type.}
\label{tab:f31_rcmax_benchmark}
{\scriptsize\setlength{\tabcolsep}{3.5pt}
\begin{tabular}{@{}lrrrr@{}}
\toprule
Comparison & W/L/T & $\bar d$ (\%) or $\rho$ & 90\% interval & $p_{\mathrm{Holm}}$ \\
\midrule
B1: IG+LS vs.\ ILS at equal CPU (mechanism) & 529/99/72 & -2.135 & [-2.347, -1.926] & $6\times10^{-68}$ \\
B2a: \textsc{QL} vs.\ random selection, equal $T$ & 229/213/258 & -0.055 & [-0.106, -0.004] & 0.150 \\
B2b: \textsc{QL} vs.\ random selection, equal CPU & 218/222/260 & -0.033 & [-0.084, +0.016] & 0.483 \\
B3: Spearman(headroom, policy gain), 28 cells & -- & +0.74 & -- & $1\times10^{-5}$ \\
B4: validated strategy (per type) vs.\ \textsc{QL} & 249/182/269 & -0.066 & [-0.120, -0.015] & 0.007 \\
\midrule
\multicolumn{5}{l}{Mean GAP (\%) to the best known solution: IG+LS \textsc{QL} / ILS at equal CPU (type)} \\
\quad job-correlated & \multicolumn{4}{l}{1.79 / 1.84} \\
\quad machine-correlated & \multicolumn{4}{l}{1.87 / 6.60} \\
\quad $U[1000,1100]$ & \multicolumn{4}{l}{0.13 / 0.25} \\
\quad $U[100,120]$ & \multicolumn{4}{l}{0.33 / 0.93} \\
\quad $U[100,200]$ & \multicolumn{4}{l}{0.66 / 2.26} \\
\quad $U[10,100]$ & \multicolumn{4}{l}{4.17 / 7.62} \\
\quad $U[1,100]$ & \multicolumn{4}{l}{5.64 / 12.34} \\
\bottomrule
\end{tabular}}
\end{table}


\subsubsection*{A prospective test of the headroom law}
The law was derived from, and tested on, destruction strategies. To test it prospectively in an action space with headroom, an IG started from the rounded LP relaxation (LP-IG, described below) was run with the
destruction size $d/n\in\{0.01,0.02,0.05,0.10,0.20\}$ as the action on the $700$ benchmark instances, with the design fixed beforehand (Table~\ref{tab:f33_dpolicy}), against random choice of $d$, the choice validated on $10$ separate instances
per type and job count, online Q-Learning, UCB1 and a table pre-trained offline. The headroom is real but modest: validation improves on random choice by $0.065\%$ per $n$ and $0.201\%$ per type and $n$ (per-instance
oracle: $0.325\%$). The policy does not capture it. Online Q-Learning is \emph{worse} than random choice ($+0.152\%$, $p_{\mathrm{Holm}}=10^{-5}$) and than the validated choice ($+0.360\%$), the pre-trained table also loses to validation ($+0.249\%$),
and over the $28$ cells the policy gain is unrelated to the headroom (Spearman $-0.37$; the criterion fixed beforehand was $\geq0.5$; Figure~\ref{fig:headroom}); the largest loss ($+0.44\%$) is in $U[1,100]$, where the headroom is largest. Most of the predictions
recorded before the run were wrong. \textbf{Headroom is therefore necessary but not sufficient.} An exploratory diagnostic suggests why: the immediate reward ranks the sizes consistently with their final quality, but more than $90\%$ of it arrives in
the first $100$ of the $1{,}000$ iterations, leaving the agent a very short window to learn; and bandits designed for such a window (Thompson sampling, with and without discounting, and a successive-elimination race) did not beat random choice either in a
pilot on $126$ tuning instances, so none of five learners captures the headroom. Choosing $d$ is, in these data, a configuration problem solved offline, not a control problem solved online.

\paragraph{Learners of other classes}
Since all five learners above are tabular or multi-armed, four learners of other classes were tested on the same $700$ instances (two in this paragraph, two neural ones in the next), with the design and the predictions recorded beforehand (Table~\ref{tab:f37_ctx}). A contextual bandit with linear function approximation (LinUCB; context: progress, stagnation and two solver features;
$\alpha$ chosen on tuning instances) gains $0.039\%$ over random choice of $d$ (interval $-0.092$ to $+0.013$; $p_{\mathrm{Holm}}=0.69$; $19\%$ of the validated gain) and loses $0.159\%$ to the validated choice ($p_{\mathrm{Holm}}=9\times10^{-15}$); its cost is $1.03$ times that of random choice. A selector trained \emph{offline} on the $280$ tuning instances from eight features of the time matrix, without the instance type
(a nearest-neighbour regression chosen by cross-validation), does capture half of the validated gain ($0.100\%$ over random choice; $p_{\mathrm{Holm}}=5\times10^{-4}$) but is worse than validation per type and $n$ by $0.105\%$ ($p_{\mathrm{Holm}}=3\times10^{-9}$; $38$ wins against $116$ losses); its cross-validated ratio on the tuning instances already ranked it behind the validated choice. Neither changes the conclusion: the class of the learner is not what prevents the capture of the headroom,
and with $280$ tuning instances the choice validated per cell remains the best configuration found. The gain of the instance-conditional selector does follow the headroom across the $28$ cells (Spearman $0.64$), as the law predicts for a selector trained offline.

% Auto-generated by analysis/make_table_ctx.py (Entrada 159/160)
\begin{table}[t]
\centering
\caption{Learners of other classes for the destruction size ($\mathrm{R}||C_{\max}$, standard benchmark, LP-IG, the $700$ test instances of Table~\ref{tab:f33_dpolicy}, $T=1000$; design and predictions recorded beforehand). LinUCB: contextual bandit with linear function approximation (context: progress, stagnation and two solver features; $\alpha=0.25$ chosen on $280$ tuning instances). Instance-conditional selector: $k$-nearest-neighbour regression on eight features of the time matrix (no instance type), trained offline on the same $280$ tuning instances and chosen by cross-validation among three models. $\bar d$ in \% of the makespan (negative: first better); Wilcoxon, Holm over four; confirmed: $p_{\mathrm{Holm}}<0.05$ and $\bar d<0$. Last row: gain over random $d$ as a fraction of the gain of the validated choice ($0.201\%$).}
\label{tab:f37_ctx}
{\scriptsize\setlength{\tabcolsep}{3pt}
\begin{tabular}{@{}lrrrrl@{}}
\toprule
Comparison & W/L/T & $\bar d$ (\%) & 90\% interval & $p_{\mathrm{Holm}}$ & Confirmed \\
\midrule
C1: offline selector vs validated per cell & 38/116/546 & +0.105 & [+0.056, +0.158] & $3\times10^{-9}$ & no \\
C2: offline selector vs random $d$ & 196/151/353 & -0.100 & [-0.158, -0.043] & $5\times10^{-4}$ & yes \\
C3: LinUCB vs random $d$ & 164/168/368 & -0.039 & [-0.092, +0.013] & 0.69 & no \\
C4: validated per cell vs LinUCB & 234/91/375 & -0.159 & [-0.218, -0.104] & $9\times10^{-15}$ & yes \\
\midrule
\multicolumn{6}{p{0.95\textwidth}}{Capture of the validated gain: instance-conditional selector 0.50; LinUCB 0.19; UCB1 0.09; online Q-Learning -0.76; pre-trained Q-Learning -0.20} \\
\bottomrule
\end{tabular}}
\end{table}


Two neural learners were added with the same design (Table~\ref{tab:f39_neural}). A one-hidden-layer Q-network trained online by temporal differences (learning rate chosen on the tuning instances) gains $0.005\%$ over random choice of $d$ (interval $-0.061$ to $+0.046$;
$p_{\mathrm{Holm}}=0.17$; $3\%$ of the validated gain) at $1.5$ times the CPU time of random choice, and loses $0.192\%$ to the validated choice ($p_{\mathrm{Holm}}=6\times10^{-18}$). A multilayer perceptron trained offline on the $280$ tuning instances (same features, standardised targets) captures about half of the validated
gain ($0.099\%$ over random choice; the interval of the mean excludes zero, $-0.162$ to $-0.036$, but the Wilcoxon test does not, $p_{\mathrm{Holm}}=0.17$, contrary to the prediction recorded beforehand) and is worse than validation per type and $n$ by $0.104\%$ ($p_{\mathrm{Holm}}=8\times10^{-15}$; $49$ wins against $192$ losses).
None of the learners above beats the validated choice, and no online learner captures more than a fifth of its gain (larger networks and a policy-gradient learner, below).

% Auto-generated by analysis/make_table_neural.py (Entrada 162/166)
\begin{table}[t]
\centering
\caption{Neural learners for the destruction size (same $700$ test instances and actions as Tables~\ref{tab:f33_dpolicy} and~\ref{tab:f37_ctx}; design and predictions recorded beforehand). Neural Q-learning: a one-hidden-layer network ($16$ units) trained online by temporal differences on the context of Table~\ref{tab:f37_ctx} (learning rate chosen on $280$ tuning instances). Offline MLP selector: a two-layer perceptron ($16\times16$) trained on the same $280$ tuning instances to predict the ratio to the best known solution of each size from eight features of the time matrix (targets standardised per size). $\bar d$ in \% of the makespan (negative: first better); Wilcoxon, Holm over four; confirmed: $p_{\mathrm{Holm}}<0.05$ and $\bar d<0$. Last row: gain over random $d$ as a fraction of the gain of the validated choice.}
\label{tab:f39_neural}
{\scriptsize\setlength{\tabcolsep}{3pt}
\begin{tabular}{@{}lrrrrl@{}}
\toprule
Comparison & W/L/T & $\bar d$ (\%) & 90\% interval & $p_{\mathrm{Holm}}$ & Confirmed \\
\midrule
N1: neural Q-learning vs random $d$ & 142/190/368 & -0.005 & [-0.061, +0.046] & 0.17 & no \\
N2: validated per cell vs neural Q-learning & 271/98/331 & -0.192 & [-0.252, -0.134] & $6\times10^{-18}$ & yes \\
N3: offline MLP selector vs validated per cell & 49/192/459 & +0.104 & [+0.065, +0.145] & $8\times10^{-15}$ & no \\
N4: offline MLP selector vs random $d$ & 169/197/334 & -0.099 & [-0.162, -0.036] & 0.17 & no \\
\midrule
\multicolumn{6}{p{0.95\textwidth}}{Capture of the validated gain: offline MLP selector 0.49; neural Q-learning 0.03; tabular Q-learning -0.76; UCB1 0.09; pre-trained Q-learning -0.20} \\
\bottomrule
\end{tabular}}
\end{table}


\paragraph{More training data: the offline selector catches up}
The instance-conditional selectors were trained on $280$ instances, which invites the objection that this is too little. The $420$ reserved instances were run with the five fixed sizes and added to the training set ($700$ instances, disjoint from the test instances; design and predictions recorded beforehand; Table~\ref{tab:f41_bigtrain}).
Validation per cell hardly changes ($+0.014\%$ against the $280$-instance version, not significant); the nearest-neighbour selector gets closer ($+0.105\%\to+0.025\%$ worse than validation, interval $-0.015$ to $+0.068$) and the MLP selector \emph{matches} it: the mean difference falls from $+0.104\%$ to $-0.011\%$ (interval $-0.046$ to $+0.019$),
with a gain of $0.198\%$ over random choice against $0.187\%$ for validation per cell. The signed-rank test flags the last comparison as confirmed in favour of the MLP ($p_{\mathrm{Holm}}=10^{-8}$), but the MLP is worse in $160$ instances and better in $36$ and the mean is equivalent to zero: the reading is \textbf{parity, not superiority}, and the decision rule, which paired the rank test with the sign of the mean,
did not anticipate their disagreement (the prediction recorded beforehand was a non-confirmation with a positive difference). What matters for the attribution is where learning pays: offline, from data of other instances, and not in the first hundred iterations of one run.

% Auto-generated by analysis/make_table_bigtrain.py (Entrada 168/169)
\begin{table}[t]
\centering
\caption{More offline training data for the instance-conditional selectors (unrelated machines, destruction size; the $280$ tuning instances plus the $420$ reserved ones, $700$ in all, disjoint from the $700$ test instances; design and predictions recorded beforehand). $\bar d$ in \% of the makespan (negative: first better); Wilcoxon, Holm over four; confirmed: $p_{\mathrm{Holm}}<0.05$ and $\bar d<0$. \emph{Caution on H2:} the signed-rank test flags the MLP as confirmed, but it is worse in $160$ instances and better in only $36$ (rank-biserial $+0.76$); the mean is $-0.011\%$ with an interval that includes zero (equivalent within $\pm0.05\%$). The honest reading is parity, not superiority. Lower block (descriptive): the same selectors against validation per cell with the same training set.}
\label{tab:f41_bigtrain}
{\scriptsize\setlength{\tabcolsep}{3pt}
\begin{tabular}{@{}lrrrrl@{}}
\toprule
Comparison & W/L/T & $\bar d$ (\%) & 90\% interval & $p_{\mathrm{Holm}}$ & Confirmed \\
\midrule
kNN selector ($700$) vs validated per cell ($700$) & 32/51/617 & +0.025 & [-0.015, +0.068] & 0.24 & no \\
MLP selector ($700$) vs validated per cell ($700$) & 36/160/504 & -0.011 & [-0.046, +0.019] & $1\times10^{-8}$ & yes \\
validated per cell: $700$ vs $280$ training instances & 46/22/632 & +0.014 & [-0.004, +0.035] & 0.43 & no \\
validated per cell ($700$) vs random $d$ & 249/90/361 & -0.187 & [-0.248, -0.131] & $1\times10^{-16}$ & yes \\
\midrule
\quad kNN, $280$ training instances vs validated per cell & 38/116/546 & +0.105 & [+0.056, +0.158] & -- & -- \\
\quad kNN, $700$ training instances vs validated per cell & 32/51/617 & +0.025 & [-0.015, +0.068] & -- & -- \\
\quad MLP, $280$ training instances vs validated per cell & 49/192/459 & +0.104 & [+0.065, +0.145] & -- & -- \\
\quad MLP, $700$ training instances vs validated per cell & 36/160/504 & -0.011 & [-0.046, +0.019] & -- & -- \\
\bottomrule
\end{tabular}}
\end{table}


\paragraph{A larger neural policy and a larger action space}
The one-hidden-layer network above is small, and the action spaces so far hold five arms or fewer. A DQN with two hidden layers of $64$ units, a replay buffer, a target network and a context that includes the success rate of each action was run, under a design and predictions recorded beforehand (Table~\ref{tab:f42_dqn}),
in three environments: the destruction size of the LP-started IG on the standard unrelated-machine benchmark ($700$ test instances), the destruction size of the IG for identical machines ($12$ classes, $1{,}200$ instances) and a new joint space of $15$ actions, three destruction strategies by five sizes, in the same identical-machine instances
(the validated choice is the best of the $15$ fixed actions per class, chosen on separate instances). Only the learning rate was tuned. In no environment does the DQN beat random choice of the same actions: it gains $0.025\%$ for unrelated machines (interval $-0.100$ to $+0.047$; $12\%$ of the validated gain) and loses $0.006\%$ and $0.023\%$
for identical machines (the latter detectable by the rank test, $p_{\mathrm{Holm}}=0.006$, but within $\pm0.05\%$ of zero); validation beats it by $0.169\%$, $0.094\%$ and $0.094\%$ ($p_{\mathrm{Holm}}\leq2\times10^{-16}$). The size of the network does not matter: against the small network the difference is $-0.018\%$ ($p=0.54$).
In the joint space the learners of three classes (tabular Q-Learning, LinUCB, DQN) are level with or slightly worse than random choice ($+0.013\%$, $+0.004\%$, $+0.023\%$) and validation beats each by $0.075$--$0.094\%$. The headroom of the larger space ($0.11\%$ for the per-instance oracle, against $0.13\%$ for the five sizes; the prediction recorded beforehand was at least $0.25\%$) is small, and across cells it does not
predict the gain of the DQN ($\rho=+0.16$, $-0.43$ and $-0.52$ in the three environments): the law fails again for the destruction size, as it did prospectively for the tabular agent. The DQN costs $0.8$, $1.8$ and $2.6$ times the CPU time of random choice. The reading is that of the previous paragraphs: the failure of the online policy does not depend on the size of the learner or, within this family, on the size of the action space;
sixteen of the eighteen recorded predictions held (the misses: the headroom of the joint space and the CPU ratio of the DQN in it).

% Auto-generated by analysis/make_table_dqn.py (Entrada 175)
\begin{table}[t]
\centering
\caption{A larger neural policy and a larger action space do not change the conclusion. Online DQN ($2\times64$ units, replay, target network, success history in the context; learning rate chosen on tuning instances) against uniform random choice of the same actions (I1: negative favours the DQN) and against the configuration validated offline on separate instances (I2: negative favours validation); $\bar d$ in \% of the makespan with the 90\% interval, Holm over six tests. Oracle: best fixed action per instance against random choice (headroom); capture: gain of the DQN over random choice divided by the gain of the validated choice (negative: worse than random); $\rho$: Spearman correlation across cells between the headroom and the gain of the DQN; CPU: ratio to random choice.}
\label{tab:f42_dqn}
{\scriptsize\setlength{\tabcolsep}{2.5pt}
\begin{tabular}{@{}lrlrlrrrrr@{}}
\toprule
Environment (actions) & $N$ & I1: DQN vs random & $p_{\mathrm{H}}$ & I2: validated vs DQN & $p_{\mathrm{H}}$ & Oracle & Capture & $\rho$ & CPU \\
\midrule
R$||C_{\max}$ benchmark, $d/n$ (5 actions) & 700 & -0.025 [-0.100, +0.047] & $0.53$ & -0.169 [-0.234, -0.109] & $1\times10^{-16}$ & -0.325 & +0.12 & +0.16 & 0.8 \\
P$||C_{\max}$, $d/n$ (5 actions) & 1{,}200 & +0.006 [-0.006, +0.018] & $0.53$ & -0.094 [-0.106, -0.083] & $2\times10^{-30}$ & -0.127 & -0.07 & -0.43 & 1.8 \\
P$||C_{\max}$, strategy $\times$ size (15 actions) & 1{,}200 & +0.023 [+0.012, +0.034] & $6\times10^{-3}$ & -0.094 [-0.105, -0.082] & $1\times10^{-33}$ & -0.112 & -0.32 & -0.52 & 2.6 \\
\bottomrule
\end{tabular}}
\end{table}


\paragraph{Qualitatively different actions: the operators of the ILS}
Destruction strategies and sizes resemble each other, so the same learners were run where the actions differ in effect and in cost: the four neighbourhood operators of the ILS for identical machines (inter-machine swap, reinsertion, intra-machine permutation and block move; $12$ classes, $1{,}200$ instances; design and predictions recorded beforehand; vectorized operators whose results are identical to the originals).
Here no single operator is a good static configuration: the mixture of the four chosen at random beats each of them per class, the per-instance oracle among the single operators is $0.04\%$ worse than the mixture, and the single operator validated per class (always the inter-machine swap) loses to the mixture by $0.062\%$. The DQN and LinUCB lose to the mixture too, by $0.041\%$ and $0.066\%$ ($p_{\mathrm{Holm}}=4\times10^{-4}$ and $1\times10^{-6}$),
and the tabular agent by $0.044\%$ on the $780$ instances stored for it; the DQN is level with the validated single operator ($-0.021\%$, $p=0.063$). Two remarks follow. The competitor that wins in this space is a static mixture, not a validated single action, so the offline-versus-online rule of this article holds if ``static'' includes mixtures. And with the efficient implementation the agents are not cheap:
the DQN costs $3.3$ times and LinUCB $1.6$ times the CPU time of random selection (the prediction recorded beforehand, $1.0$--$1.6$ and $0.9$--$1.3$, missed), so at equal CPU they would fare worse; the other four predictions held.

\paragraph{Compiled agents, more repetitions and spaces of up to 30 actions}
The agents above run in Python and the policy effect is small, so it was measured again with a compiled IG and three compiled agents: tabular Q-Learning, UCB1 and a policy-gradient learner (a softmax policy per state with a baseline; learning rate chosen on tuning instances). Compilation makes a run about a hundred times faster, which allows five seeds per instance (the unit of analysis is the mean of the five)
in four identical-machine action spaces: three destruction strategies, five sizes, their $15$ combinations and $30$ actions that also choose whether to apply the local search (design and predictions recorded beforehand; Table~\ref{tab:f43_agents}). The compiled agents reproduce the Python ones (differences from $-0.003\%$ to $+0.012\%$ in five checks, all within $\pm0.05\%$). A compiled agent is cheap per unit of work, $1.0$--$1.09$ times for UCB1 and the policy gradient
and $1.07$--$1.33$ times for Q-Learning, which nevertheless needs only $0.50$--$0.60$ of the CPU time of random choice in the three larger spaces because it learns to pick the cheaper actions ($0.38$--$0.50$ of the evaluations). At an equal number of iterations no learner gains more than $0.016\%$ over random choice in any of the $12$ comparisons, within the equivalence margin
(Q-Learning loses by $0.017\%$ in the $30$-action space). At equal CPU time Q-Learning does win in the two spaces of sizes, by $0.059\%$ and $0.049\%$ (both beyond the margin), but only by choosing cheaper actions, since at equal iterations its gain is nil. The configuration validated per class beats every learner by $0.076$--$0.095\%$ in the spaces of $5$ and $15$ actions ($p_{\mathrm{Holm}}<10^{-58}$) and is level with them
(within $0.03\%$) in the other two, where the headroom of the per-instance oracle is smaller ($0.05\%$ against $0.12$--$0.13\%$). Across the $48$ cells (class by space) the headroom does not predict the gain of Q-Learning ($\rho=0.09$). Sixteen of the nineteen recorded predictions held (the misses: the equal-CPU verdict of Q-Learning in the strategy space, its size in the $30$-action space, and the headroom of the latter, predicted at $0.1$--$0.6\%$).

% Auto-generated by analysis/make_table_agents.py (Entrada 181)
\begin{table}[t]
\centering
\caption{Compiled agents in action spaces of $3$ to $30$ actions (identical machines, $12$ classes $\times$ $100$ instances, $5$ seeds per instance, the unit of analysis is the mean of the five; $1{,}200$ pairs per row). Oracle: gain of the best fixed action per instance over random choice (\%). $K$: CPU time of the learner over that of random choice. Equal iterations: learner vs random choice ($T=1000$). Equal CPU: learner vs random choice with $T=\mathrm{round}(1000K)$ iterations. Validated: the best fixed action per class (chosen on $30$ separate instances per class) vs the learner, $T=1000$ (negative: the validated choice wins). $\bar d$ in \% of the makespan with the 90\% interval; Holm over $12$ tests in each of the three comparisons.}
\label{tab:f43_agents}
{\scriptsize\setlength{\tabcolsep}{2.5pt}
\begin{tabular}{@{}lllrlll@{}}
\toprule
Actions & Oracle & Learner & $K$ & Equal iterations & Equal CPU & Validated vs learner \\
\midrule
\multirow{3}{*}{3 strategies} & \multirow{3}{*}{$0.046$} & Q-Learning & 1.07 & -0.016 [-0.021, -0.011] & -0.011 [-0.016, -0.006] & +0.016 [+0.010, +0.023] \\
 &  & UCB1 & 1.01 & -0.002 [-0.007, +0.003] & +0.000 [-0.005, +0.005] & +0.002 [-0.004, +0.009] \\
 &  & Policy gradient & 1.01 & +0.003 [-0.002, +0.008] & +0.005 [+0.000, +0.010] & -0.003 [-0.009, +0.003] \\ \midrule
\multirow{3}{*}{5 sizes} & \multirow{3}{*}{$0.125$} & Q-Learning & 0.60 & +0.003 [-0.005, +0.011] & -0.059 [-0.068, -0.050] & -0.095 [-0.102, -0.087] \\
 &  & UCB1 & 1.01 & -0.001 [-0.006, +0.005] & +0.006 [+0.001, +0.011] & -0.091 [-0.097, -0.085] \\
 &  & Policy gradient & 1.03 & +0.000 [-0.005, +0.005] & +0.005 [-0.000, +0.010] & -0.092 [-0.098, -0.086] \\ \midrule
\multirow{3}{*}{15 (strategy $\times$ size)} & \multirow{3}{*}{$0.122$} & Q-Learning & 0.59 & +0.004 [-0.004, +0.012] & -0.049 [-0.058, -0.040] & -0.084 [-0.092, -0.076] \\
 &  & UCB1 & 1.03 & -0.004 [-0.009, +0.001] & -0.001 [-0.005, +0.004] & -0.076 [-0.082, -0.070] \\
 &  & Policy gradient & 1.06 & -0.000 [-0.005, +0.005] & +0.008 [+0.003, +0.013] & -0.080 [-0.086, -0.074] \\ \midrule
\multirow{3}{*}{30 (+ local search)} & \multirow{3}{*}{$0.049$} & Q-Learning & 0.50 & +0.017 [+0.011, +0.023] & -0.016 [-0.022, -0.010] & -0.029 [-0.035, -0.022] \\
 &  & UCB1 & 1.05 & -0.004 [-0.007, -0.000] & -0.002 [-0.005, +0.001] & -0.008 [-0.012, -0.004] \\
 &  & Policy gradient & 1.08 & +0.001 [-0.003, +0.005] & +0.002 [-0.001, +0.006] & -0.013 [-0.017, -0.008] \\
\bottomrule
\end{tabular}}
\end{table}


\begin{figure}[t]
\centering
\includegraphics[width=0.95\textwidth]{fig_headroom.pdf}
\caption{Headroom of the actions against the gain of the Q-Learning policy over random selection, per cell (type $\times$ $n$; $28$ cells) of the standard $\mathrm{R}||C_{\max}$ benchmark, $700$ test instances. Headroom is the gain of the
per-instance oracle over random selection (negative where selection noise outweighs the differences between actions). A: destruction strategies, where the law holds (Spearman $\rho=0.74$). B: destruction size, a prospective test
with $\rho=-0.37$: the largest headroom goes with the largest losses of the policy (below zero).}
\label{fig:headroom}
\end{figure}

% Auto-generated by analysis/make_tables_cycle5.py (Entrada 129)
\begin{table}[t]
\centering
\caption{A prospective test of the headroom law with the destruction size as the action (LP-IG, 700 instances of the standard $\mathrm{R}||C_{\max}$ benchmark, $T=1000$; actions $d/n\in\{0.01,0.02,0.05,0.10,0.20\}$). Validated choices and the pre-trained tables use 10 tuning instances per type and $n$. $\bar d$ in \% of the makespan (negative: first better); Wilcoxon, Holm over five. D5 is the Spearman correlation between the headroom of a cell (type $\times$ $n$; 28 cells) and the policy gain, one-sided, with the pre-specified criterion $\rho\geq0.5$. Last row: mean gain over random $d/n$ (\%) of each choice.}
\label{tab:f33_dpolicy}
{\scriptsize\setlength{\tabcolsep}{3.5pt}
\begin{tabular}{@{}lrrrr@{}}
\toprule
Comparison & W/L/T & $\bar d$ (\%) or $\rho$ & 90\% interval & $p_{\mathrm{Holm}}$ \\
\midrule
D1: validated $d/n$ (per $n$) vs.\ random $d/n$ & 189/181/330 & -0.065 & [-0.126, +0.004] & 0.006 \\
D2: \textsc{QL} over $d/n$ vs.\ random $d/n$ & 144/221/335 & +0.152 & [+0.080, +0.223] & $1\times10^{-5}$ \\
D3: \textsc{QL} vs.\ validated $d/n$ (per type, $n$) & 85/298/317 & +0.360 & [+0.285, +0.443] & $2\times10^{-30}$ \\
D4: pre-trained \textsc{QL} vs.\ validated $d/n$ (per type, $n$) & 96/234/370 & +0.249 & [+0.184, +0.317] & $2\times10^{-16}$ \\
D5: Spearman(headroom, \textsc{QL} gain), 28 cells & -- & -0.37 & -- & 0.975 \\
\midrule
\multicolumn{5}{p{0.96\textwidth}}{Gain over random $d/n$ (\%): validated per $n$ -0.065; validated per type and $n$ -0.201; per-instance oracle -0.325; \textsc{QL} +0.152; UCB1 -0.018; pre-trained \textsc{QL} +0.040} \\
\bottomrule
\end{tabular}}
\end{table}


\subsubsection*{A stronger learning-free heuristic: LP start and a restricted MILP}
The distance to the exact references suggested improving the heuristic rather than the policy. The linear relaxation of the assignment model (solved in $0.4$ seconds on average) has a vertex solution with
fewer fractional jobs than machines ($29$ of $1{,}000$ jobs in the largest cells); rounding it (integral jobs stay, fractional ones go to the machine of earliest completion within their support) and descending gives a start that is
poor by itself ($5.2\%$ above the best known solution) but from which the same IG is far better. On the $420$ reserved instances, with $d/n$ tuned on $280$ others and a pre-specified design, this LP-IG beats the IG started
from the constructive solution by $1.18\%$ at equal CPU time ($308$ wins against $21$), and the mean GAP to the best known solutions falls from $2.23\%$ to $0.97\%$. A MILP restricted to the six cheapest machines of each job (and to the
machine of the incumbent, so that the restricted problem contains the LP-IG solution), solved with HiGHS for at most $30$ seconds, never worsens the solution and improves $54.8\%$ of them; against more iterations of the LP-IG at a
similar CPU time ($T=8000$) it is better by $0.49\%$ ($196$ wins against $33$, $p_{\mathrm{Holm}}=10^{-24}$), and the GAP falls to $0.25\%$ ($0.11\%$ for $n=1000$; Table~\ref{tab:f34_lpigsr}). That is within the range of CPLEX with two hours
($0.20\%$) and $0.15$ points from Gurobi and the exact method with $300$ seconds. The criterion fixed beforehand ($\leq0.3\%$ overall and $\leq0.5\%$ in every job count and type) is met overall and for every $n$ but not for
job-correlated instances ($0.67\%$). A variant that chooses the candidate machines by the reduced cost of the LP was tested on the remaining $700$ instances: significant in sign but negligible in size ($-0.03\%$;
mean GAP $0.30\%\to0.27\%$), and its criterion ($\leq0.25\%$) was not met. No learned component takes part in any of the three steps.

% Auto-generated by analysis/make_tables_cycle5.py (Entrada 129)
\begin{table}[t]
\centering
\caption{LP-IG followed by a MILP restricted to the six cheapest machines of each job (and to the machine of the incumbent), 30\,s limit, on the 420 reserved instances of the standard $\mathrm{R}||C_{\max}$ benchmark (design on 168 tuning instances; $T=1000$). $\bar d$ in \% of the makespan (negative: first better); Wilcoxon, Holm over three. Below the line, the mean GAP (\%) to the best known solution by $n$: baseline IG / LP-IG ($T=1000$) / LP-IG ($T=8000$) / LP-IG + restricted MILP / Gurobi / LR (CPLEX in 2\,h: 0.20). Mean CPU time per run: LP-IG 3.3\,s; LP-IG ($T=8000$) 27.3\,s; + restricted MILP 23.3\,s.}
\label{tab:f34_lpigsr}
{\scriptsize\setlength{\tabcolsep}{3.5pt}
\begin{tabular}{@{}lrrrr@{}}
\toprule
Comparison & W/L/T & $\bar d$ (\%) & 90\% interval & $p_{\mathrm{Holm}}$ \\
\midrule
S1: + restricted MILP vs.\ LP-IG with $T=8000$ (similar CPU) & 196/33/191 & -0.489 & [-0.628, -0.366] & $6\times10^{-25}$ \\
S2: + restricted MILP vs.\ LP-IG ($T=1000$) & 230/0/190 & -0.675 & [-0.832, -0.540] & $5\times10^{-39}$ \\
S3: + restricted MILP vs.\ LP-IG ($T=5000$) & 204/28/188 & -0.552 & [-0.709, -0.420] & $2\times10^{-27}$ \\
\midrule
\multicolumn{5}{l}{Mean GAP (\%) to the best known solution} \\
\quad $n=100$ & \multicolumn{4}{l}{1.85 / 1.21 / 0.77 / 0.36 / 0.26 / 0.00} \\
\quad $n=200$ & \multicolumn{4}{l}{2.48 / 1.74 / 1.43 / 0.27 / 0.09 / 0.01} \\
\quad $n=500$ & \multicolumn{4}{l}{2.41 / 0.63 / 0.59 / 0.28 / 0.04 / 0.20} \\
\quad $n=1000$ & \multicolumn{4}{l}{2.18 / 0.31 / 0.31 / 0.11 / 0.02 / 0.15} \\
\quad all instances & \multicolumn{4}{l}{2.23 / 0.97 / 0.78 / 0.25 / 0.10 / 0.09} \\
\bottomrule
\end{tabular}}
\end{table}

